\centerline{\bf \Huge Уголки и дуги}

\begin{flushright}
        \scriptsize{Эпизод 1. Нападение Ангела.}
\end{flushright}

\problem{0}
Оптимизационные леммы. На окружности выбраны точки $A, B, C$ и $D$ (именно в таком порядке по часовой стрелке)

(a) Отрезки $A C$ и $B D$ пересекаются в точке $P$. Докажите, что $\angle A P D=\frac{1}{2}(\overparen{A D}+\overparen{C B})$

(b) Лучи $A B$ и $D C$ пересекаются в точке $Q$. Докажите, что $\angle A Q D=\frac{1}{2}(\overparen{A D}-\overparen{C B})$

(c) Докажите, что $A D \| B C$ тогда и только тогда, когда $\overparen{A B}=\overparen{C D}$

\problem{1}
Шестиугольник $A B C D E F$ вписан в окружность. Докажите, что

$$
\angle A+\angle C+\angle E=\angle B+\angle D+\angle F
$$

\problem{2}
Треугольник $A B C$ вписан в окружность $\omega$. Биссектрисы его внутренних углов $B$ и $C$ пересекают $\omega$ в точках $X$ и $Y$. Докажите, что прямая $X Y$ отсекает от угла $B A C$ равнобедренный треугольник.

\problem{3} 
Шестиугольник $A_1 A_2 A_3 A_4 A_5 A_6$ вписан в окружность. Оказалось, что $A_1 A_2 \| A_4 A_5$ и $A_2 A_3 \| A_5 A_6$. Докажите, что $A_3 A_4 \| A_6 A_1$.

\problem{4}
Точки $A$ и $B$ движутся по окружности $\omega$ так, что величина дуги $A B$ постоянна. Докажите, что прямая $A B$ касается некоторой фиксированной окружности, не меняющейся с течением времени.

\problem{5}
Четырехугольник $A B C D$ вписан в окружность. Докажите, что биссектрисы углов $A B D$ и $A C D$ пересекаются на этой окружности.

\problem{6}
В окружность вписан четырёхугольник $ABCD$. Его вершины разбивают окружность на четыре дуги. Докажите, что отрезки, соединяющие середины противоположных дуг, перпендикулярны.

\problem{7}
Дан выпуклый четырехугольник $ABCD$, в котором $AB=BC=CD$. Его диагонали $A C$ и $BD$ пересекаются в точке $E$. Описанная окружность треугольника $A D E$ пересекает сторону $A B$ в точке $P \neq A$ и продолжение стороны $C D$ в точке $Q \neq D$. Докажите, что отрезки $A P$ и $D Q$ равны.

\problem{8}
Точка $P$ движется по «меньшей» дуге $B C$ описанной окружности треугольника $A B C$ (треугольник фиксирован). Точка $Q$ на отрезке $B C$ выбирается так, что $\angle B P Q=\angle A B C$. Докажите, что все прямые $P Q$ имеют общую точку.

\problem{9}
Три равные окружности имеют общую точку $H$ и попарно повторно пересекаются в точках $A$, $B$ и $C$, причем точка $H$ лежит внутри треугольника $A B C$. Докажите, что $H$ - ортоцентр (то есть, точка пересечения высот) треугольника $A B C$.